\documentclass[11pt,a4paper]{article}
\usepackage[T1]{fontenc}
\usepackage{lmodern}
\usepackage[margin=27mm]{geometry}
\usepackage{amsmath,amssymb,microtype}
\usepackage[hidelinks]{hyperref}
\usepackage{enumitem}
\setlist{itemsep=3pt,topsep=5pt}
\emergencystretch=2em
\title{Lean Verification Plan for Dense Max 2-SAT}
\author{}
\date{23 October 2026}
\begin{document}
\maketitle

The active manuscript is \texttt{dense\_max2sat\_arxiv.tex}. The aim is
for every retained formal result to have a corresponding Lean proof.
This is a long-term target, not a claim about current coverage. We keep
executable algorithms only when they implement procedures presented in
the paper. Supporting definitions and technical proof lemmas remain in
Lean without requiring a numbered paper statement for each one.

\section{Keep the paper and the formal artifact aligned}

For each numbered result or proposed theorem, record its label, hypotheses, conclusion,
corresponding Lean declarations, and verification status. The record is
\texttt{lean/verification/paper-correspondence.json}. The checking script
requires an entry for every numbered result and audits the named Lean
declarations. It records the manuscript and source hashes. It also flags changes to a statement for renewed correspondence review.
Mathematical correspondence still requires review,
especially when a statement changes without changing its label.

There are four statuses: checked mathematical content; checked core with
an ordinary arithmetic-cost argument; partial mathematical verification;
and unverified. A theorem with unfinished internal proof obligations is
not marked complete. Published dependencies are permitted under the policy
below, with the conditional guarantee explicitly recorded. Correctness, candidate counting, recursion-node
bounds, arithmetic cost, and storage bounds are separate obligations.

Every paper revision affecting a formalized result must update its Lean
statement or explicitly record a remaining discrepancy. Rebuild and audit
before reporting a new verification result. Earlier artifact archives
remain historical snapshots; they are not additional active algorithms.

\section{Completed mathematical foundation}

The first milestone supplies clause scores, score-preserving unit and pair
reductions, executable preprocessing, exact strict-improvement constraints,
propagation, complementary cancellation, and sound summed-constraint
rejection. It also proves guarded domination and exact output-sensitive
enumeration of a single nonnegative rational inequality. Passing the sum
check does not certify feasibility of a collection of inequalities.

The formal input permits clause multiplicities, with distinct variables
inside each binary clause. This covers simple original inputs and the
repeated units created by preprocessing. Canonical-form counting facts
needed by later distributional results must still be proved where used.

\section{Completed literal-gap mathematical bounds}

The paper now uses the verified fixed-preference algorithm: preprocess,
fix preferences from the effective gaps, choose the half carrying more
total absolute gap, branch on that half, and enumerate the suffixes
satisfying each prefix's summed constraint. The older adaptive algorithm
is outside the active build.

The checked proof includes the common optimal assignment, the expected
comparison statistic, the exact occurrence influence bound, a finite
Boolean bounded-differences inequality, and the general count
\[
 W\le2^n\exp(-G^2/(512Q)),\qquad
 G=\sum_i|g_i|,\quad Q=\sum_i(D_i+u_i)^2>0.
\]
It then establishes correctness after preprocessing and the explicit
relative-bias search bound
\[
 (4+2n)2^{(1-\varepsilon^2\eta^2/(3000R))n}.
\]
Recursion depth, constraint sizes, and the agreement of streaming folds
with mathematical output lists are also checked. The polynomial
bit-overhead and space conclusion uses the paper's conventional
implementation argument. No compiler-level or physical-memory theorem
is claimed.

The dense-gap corollary is now checked as well. The formalization defines
simplicity using clause literal sets, so reversing a binary clause cannot
hide a duplicate. It proves the bound $B_i\le4(n-1)+2$, preservation of
every effective gap under preprocessing, and nonincrease of occurrence
counts. These facts supply the hypotheses of the general counting lemma
from the original input and yield the explicit bound
\[
 (4+2n)2^{(1-\varepsilon^2\eta^2/48000)n}.
\]
The arithmetic-cost verification boundary remains as stated above.

The proposed biased-subset extension is recorded without a proof in the
paper's appendix. It is not an established result and is outside the
current verification milestone. Its recognition questions and concrete
applications should not delay verification of the selected main results.

\section{Random signs: correctness, conditioning, and exposure completed}

The deterministic correctness milestone is complete. The formal procedure
computes the observed support and rational degree weights from the formula,
optimizes unit-only variables, preprocesses the core, recomputes preferences
from residual formulas, and performs weighted prefix--suffix search.
The earliest-optimum argument proves simultaneous survival of the necessary
cuts. Exact rational-sum feasibility and streaming enumeration are checked.
The final theorem returns an optimum original score and assignment, including
for repeated clauses, empty support, and zero variables.

The finite sampling models are now defined, with their symmetries and
endpoint second moments proved. The observed-support conditioning lemma
is checked, including the connection between sampled pair counts and the
neighbor relation extracted from the generated formula. The complete
exposure lemma is now checked: its auxiliary experiment has the prescribed
joint law, the gains and residual preferences are connected to the actual
algorithm, and the history determines the coefficients while leaving the
required signs fresh. The conditional mean, variance, observed-degree
bound, and later-deviation identity are formalized. The identification of
$\Delta_t$ with the average over uniform unassigned values is checked
separately from the expectation over random inputs.

The absolute-gain and second-moment lemma is now checked, including
the exact fourth moment, truncation, arbitrary shifts, and the
second-moment inequality on general probability spaces. Its proof uses
two applications of the Cauchy--Schwarz inequality, with the same bounds
as before. The finite sampling laws are connected to probability
measures by checked expectation identities.

The conditional-range lemma is now checked on finite probability spaces,
including adaptive interval lengths, a pathwise sum of squared lengths,
and the independent-input bounded-differences corollary. The zero-influence
case is included. An adaptive regression example checks that the proof
uses the budget spent on each path, rather than separate worst-case
bounds for all steps.

The accumulated-requirement and compensation lemmas are now checked in
full. Their proofs instantiate the entire fixed order and average over
the base-pattern law. For requirements, the direct pointwise estimate
$\sqrt{Y_t}\ge Y_t/(m_0\sqrt{d_t})$ replaces the variance calculation,
while retaining the theorem's constant. The independent-edge law,
coordinate sensitivities, and their summed budget are checked. For
compensation, grouping by the later endpoint gives checked conditional
exponential estimates and a total budget $m_0^2T$.

The identity between the summed comparison and $A_T+Q_T$, the joint
feasibility bound, and its transfer to the original input/assignment law
are checked. Expected and high-probability candidate counts are also
checked for a fixed support.

The node bound for the executable weighted prefix/suffix search is now
checked, including rejected calls and its formula-only initialization.
Every output satisfies the comparisons on its full prefix; the converse
is unnecessary because early pruning may remove further candidates.
The score interpretation of the actual rational constraints, preprocessing,
and the prescribed order are connected to the analytic count. This yields
expected and high-probability node bounds on a fixed support whose indexed
vertices are all active, with a fixed exponential base below two. The same
bounds hold for the conditional retained-pair law.

The arbitrary-support runtime assembly is now complete. It avoids
reindexing: original labels are retained, inactive variables spend no
influence, and inactive coordinates are fixed when candidates are counted.
Deleting empty pair draws preserves the score comparisons and candidate
count. The conditional retained-pair law, the empty-support case, the
comparison $r\le n$, and averaging over observed supports are all checked.
Lean proves expected nodes at most $(4+6n)2^ne^{-c_0n}$ and failure
probability at most $3e^{-c_0n/2}$ for the node threshold
$(4+2n)2^ne^{-c_0n/2}$. Both concrete clause models are instantiated with
the paper's common constants. No concentration or candidate bound is
assumed in the final results. The algorithm and theorem constants are
unchanged; polynomial arithmetic overhead and streamed space remain
ordinary implementation arguments.

\section{Count adjustments and graph families}

The count-adjustment algebra and optimum-survival arguments are now
formalized. This includes the actual rational score gains, the diagonal
allowances, exact counts and approximate projection tuples, preprocessing,
and the streaming prefix/suffix enumeration. The grid construction has
checked coverage, cardinality, and total-allowance bounds. The full
comparison sum is identified with accumulated gap plus compensation.
The newly stated lemma for assignments with exactly $h$ ones has direct
formal proofs of its identity and its comparison inequality for either
polarity, using the same count-adjusted gains.
The analytic accumulation proof is now complete through explicit node
bounds for both sequential comparisons and a fixed number of rounds.
The common-gap bound uses a smooth approximation to absolute value and a
finite-cube concentration inequality proved from the two-point entropy
inequality. The quadratic tail follows from a random partition and the linear
sign-sum exponential bound. Both proofs are in Lean and in the manuscript;
no concentration theorem is assumed in the final accumulation result.

The checked assembly includes the actual comparison sum, conditional block
laws, operator-norm reductions, constant substitutions, summation over all
guesses, and an exponential base strictly below two. Formula-level theorems
include preservation of count interactions by preprocessing and independence
of interaction coefficients from the guess. The main-text specialization
also checks the two-round order, prefix size, zero first-round rows, and the
$n+1$ count bound. Projection-grid corollaries combine the checked grid size
and allowance bounds with the analytic theorem. The conventional polynomial
bit-cost and space arguments are still outside the formal node-count model.

The regular-graph corollary now has a complete mathematical proof and
an explicit recursion-node bound in Lean. The proof starts from actual
adjacency eigenvalues, retaining multiplicity, and derives the centered
operator bound, mixing inequality, row energies, rounded-round profile,
and comparison allowance. It also verifies the greedy independent-set
branch, optimal unit-only completion, and their enumeration bound.
Fixed integer degree and size thresholds combine these branches with the
count search into one bound with a base below two, including arbitrary
unit clauses and small inputs. No graph estimate is left as an assumed
intermediate bound. The cited properties of Paley and Ramanujan graphs
remain external membership facts; their constructions are not formalized.
The separate random-sign runtime assembly is also complete.

The independent-pair work now proves variance-sensitive scalar tails,
squared-entry variance bounds, empirical-mean concentration, and
independent-copy moments and symmetry on finite probability spaces.
The unordered-pair sampling model is connected to the symmetric matrix,
its exact squared-energy normalization, and each earlier-neighbor row.
Deterministic centering identities and a union bound yield a simultaneous
scalar event, with an instantiation for independent Bernoulli edges.
The manuscript includes the elementary scalar proof and uses its sufficient
constants. The subsequent matrix and probability-space steps below
complete the independent-pair proposition.

The matrix estimate is now checked conditional on the explicitly quoted
Rademacher bound of Bandeira and van Handel and the convex upper-tail
inequality of Ledoux. Symmetrization, sign insertion, Euclidean norm
geometry, row and column moments, the large-density substitutions, and
the empirical perturbation are checked. The finite-law membership proof is now complete in both centering cases:
the summed row profile (including integer rounding), comparison allowance,
scalar parameter substitutions, final probability $1-1/n$, and the size
bound for count accumulation are checked. Statements for the original
coefficients identify the sampled empirical mean exactly. Bernoulli
instantiations cover homogeneous empirical centering and inhomogeneous
supplied means. The final probability-space translation also covers
independent bounded integer variables on an arbitrary probability space,
with their original expectations and measure-theoretic variances. The
proposition is stated for integer interactions, as in every formula model
in the paper; an additional claim for continuously distributed real
coefficients is outside this submission.
For hidden communities, the search wrapper now has checked exactness for
any total rational representation routine, including the small-input and
zero-energy fallbacks and the zero-coefficient branch. A concrete integer
square root computes the grid spacing; coverage, the lower cap, arbitrary-
input grid size, and inactivity of the cap on the useful event are checked.
The deterministic energy and cross-round profile perturbations are proved.
The subsequent assembly also proves the zero-coefficient runtime branch,
identifies the formula's remainder matrix, and derives the coefficient mass
and allowance from Rayleigh quotients and explicit mean/size budgets.
Cross-round scalar tails, total energy and symmetric matrix norm hold jointly
with probability at least $1-1/n$ under an explicit sufficient density bound.
The actual search consequently has a checked high-probability node bound,
conditional on the numerical representation's stated guarantees. The only
published-result hypotheses remain the two previously recorded estimates.
The manuscript uses the sufficient constant $K=400\chi$.

The next numerical batch verifies rational largest-entry sorting and the
retained representation, unit directions, exact Rayleigh coefficients,
elementary diagonal truncation, and the rank-to-Frobenius inequality.
The explicit cutoff $n\ge5{,}120{,}000r\chi$ gives the small correction
needed by the search analysis. All four previous approximation guarantees
now follow from a rational orthogonal matrix with a small diagonalization
residual. The paper presents the same elementary truncation proof; the
algorithm and input classes are unchanged.

The rational orthogonal matrix is now constructed by an executable grid
search. Exact plane rotations, the largest-pivot energy identity, a uniform
per-step decrease, and the finite iteration bound prove termination with
residual at most $n^{-4}$. Bounds in the explicit formula length show that
the grid and number of iterations are polynomial. The conventional
bit-complexity argument is recorded in the manuscript. The actual
formula-only search now has unconditional correctness and a checked
high-probability node bound under explicit model, density and size
conditions; no numerical callback or quality assumption remains.
The hidden-community correctness lemma is consequently checked.

The asymptotic parameter substitution and bounded-size running-time
assembly are now checked. The finite-law discovery theorem fixes its
rational algorithm parameter, density constant, multiplier, and exponential
base before the input size and distribution are supplied. The original
SBM specialization is also checked: independent Bernoulli edges, the actual
graph/formula encoding, mean rank and norm, minimum and maximum block
variances, and units depending arbitrarily on the sampled graph. This
closes the mathematical node bound for the main SBM theorem, conditional
on the same two recorded published estimates. Bit-operation overhead and
polynomial space retain the conventional paper argument.

The probability-space translation is now checked. Almost-sure boundedness
of integer coefficients yields a finite common alphabet. Encoding them
preserves their joint independent law, all event probabilities, moments,
and variance. Both membership alternatives and the complete discovery
family theorem apply on the original probability space. The discovery
theorem constructs a measurable event depending only on the binary
coefficients; on this event its node bound holds for every formula with
those interactions, including arbitrary subsequently chosen units.

Every numbered proved statement now has checked mathematical content;
there are no partially formalized numbered results. The independent-pair
proposition has been scoped explicitly to integer interactions, without
changing any Max 2-SAT or Max Cut input family. The only unverified numbered
statement is the deliberately unproved biased-subset research target.
Polynomial bit-operation overhead and space remain conventional paper
arguments, and the two published estimates remain explicit hypotheses.
The final audit and correspondence register document these boundaries.

\section{Trust and scope}

We formally verify our contributions and our applications of established
results, rather than undertaking to formalize every cited theorem's proof.
For each external result, record the exact published statement, source,
theorem number or page, and any rescaling or specialization. Check all
hypotheses and deductions in our application. The corresponding Lean
statement is an explicit hypothesis, never an unlabelled assumption or a
claim that the external theorem has itself been formally proved.
Verification is conditional on these named published results. The register
\texttt{lean/verification/external-results.json} distinguishes them from
unfinished lemmas specific to our paper; an unfinished internal lemma may
not be relabelled as an external result.


The project pins Lean and mathlib versions. The build rejects proof
placeholders, and the transitive axiom audit permits only Lean's standard
\texttt{propext}, \texttt{Classical.choice}, and \texttt{Quot.sound}.
Executable examples supplement universally quantified proofs.

TLA+ is not needed for the current finite combinatorial and analytic
arguments. It may later help investigate operational invariants, but
finite model checking would not replace the arbitrary-size correctness
and asymptotic proofs. We should finish paper-aligned Lean targets before
starting a second verification framework.
\end{document}
